\statement{F.13}{Proved}{The decorated two-level boundary tower has automorphism group $\GT$}
Let $\mathcal T^{\partial}_{4,5}$ consist of $\widehat T_4,\widehat T_5$, all five forgetting functors, clutching inclusions with their $\widehat T_3,\widehat T_4$ factors, and coherent relabelling pseudo-actions. Its automorphisms are continuous equivalences fixing labelled stratum types, preserving oriented boundary inertia with one exponent $\lambda\in\widehat\Z^\times$, and commuting with these operations through natural isomorphisms satisfying their coherence equations. Modulo coherent natural isomorphism,
\[
 \pi_0\Aut(\mathcal T^{\partial}_{4,5})\simeq\GT.             \tag{F.11}
\]
The same holds for the strict pro-objects. The coordinate identification uses the usual marking by oriented loops $x,y$ around $0,1$ and a compatible tangential base point. Changing the marking conjugates this description; the outer tower and arithmetic action are independent of these choices.
\begin{proof}
Restriction to the open objects is injective by F.12. At those objects the inertia and relabelling conditions are exactly the definition of $\Out_n^\#$ in \citet[definition in the introduction and Main Theorem]{HarbaterSchneps2000}. Their theorem identifies $\Out_4^\#$ with $\widehat{\mathrm{GT}}_0$ and $\Out_5^\#$ with $\GT$; forgetting gives the standard injection from the latter to the former. Thus the five-point outer automorphism determines the compatible four-point one.

Conversely their $\GT$ action preserves inertia, labelled types, and forgetting homomorphisms. F.12 extends it uniquely, up to natural isomorphism, to both boundary categories. On a divisor $d_{ij}$, forgetting $i$ or $j$ restricts to the identity of its four-point factor after relabelling; its three-point factor is trivial. Compatibility with forgetting therefore identifies the induced action with the four-point action, giving clutching compatibility. At two-node strata the component factors are all three-point categories. F.12 already forces the crossing bimodules and their compositions.

The compatibility isomorphisms are coherent. A natural automorphism of the identity on $\widehat C_n$ has open component in $Z(\widehat G_n)=1$; naturality along $j_s:s\to u$ and injectivity of postcomposition force every other component to be $1$. For forgetting, surjectivity of the open homomorphism gives the same argument, since its image has centralizer $Z(\widehat G_4)=1$. For clutching, the image is the full automorphism group of the target stratum, again centre-free or trivial. Thus compatibility isomorphisms in this two-level diagram are unique and their coherence equations hold. Relabellings are equivalences and satisfy the same uniqueness. Taking opposites preserves the result; F.2 compares the pro-object and continuous versions.
\end{proof}

\statement{F.14}{Proved}{The full genus-zero boundary tower and its arithmetic subgroup}
Let $\mathcal T^{\partial}_0$ contain all $n\geq3$ with the same decorations and all forgetting and clutching functors. Restriction is an isomorphism
\[
 \pi_0\Aut(\mathcal T^{\partial}_0)
 \xrightarrow{\sim}\pi_0\Aut(\mathcal T^{\partial}_{4,5})
 \simeq\GT.                                                \tag{F.12}
\]
F.10 gives the usual faithful arithmetic inclusion $G_{\Q}\hookrightarrow\GT$. Equating its image with all automorphisms is the open equality E.8. The categorical group thus identifies the early normalizer-type group E.7; the revised arithmetic-image group is the subgroup E.2.
\begin{proof}
The Harbater--Schneps Main Theorem and its stability corollaries identify $\Out_n^\#$ with $\GT$ for all $n\geq5$, compatibly with forgetting. Their tower action supplies existence on open strata; F.12 extends it uniquely at every level. The group of a boundary stratum is $\prod_v\widehat G_{\operatorname{val}(v)}$. Each factor is recovered by forgetting all but its markings and one marking in every branch attached to it. This restricts to the corresponding projection. Forgetting compatibility therefore forces the product of factor actions and proves clutching compatibility. All these groups have trivial centres, including the trivial valence-three factors, so F.13 proves coherence. The extension from two levels is unique.

Arithmetic descent and the usual arithmetic action coincide on the four-point open stratum by fundamental-group comparison, hence everywhere by uniqueness. Faithfulness is F.10. E.8 is the additional surjectivity assertion; no profinite equality follows from a rational Lie-algebra calculation.
\end{proof}

\subsection{F.5. Positive genus, isotropy, and the rational comparison}

\statement{F.15}{Proved}{Completion for the elliptic types retains stack isotropy}
For $(g,n)=(1,1),(1,2)$ take Jansen's labelled stable-curve category in exit variance, with a finite skeleton indexed by stable labelled dual graphs. Its strict completion F.2 is a category in profinite sets, faithful on discrete Hom-sets, and equivalent to the profinite-multicurve transporter category for $\widehat P_{g,n}$. At $\sigma$,
\[
 1\longrightarrow\prod_v\widehat P_{g_v,n_v}
 \longrightarrow\Aut_{\widehat C_{g,n}}(\sigma)
 \longrightarrow A_\sigma\longrightarrow1,                  \tag{F.13}
\]
where $A_\sigma$ is the finite image of the discrete stabilizer on the dual graph and half-edges, fixing the original marked legs. No splitting is chosen. These categories carry F.9--F.10's finite constructible-sheaf comparison and outer Galois action, now for moduli \emph{stacks}.
\begin{proof}
Jansen's Corollary 6.6 and Construction 3.18 apply in all genera, retaining graph automorphisms, and give $\Aut(\sigma)=H_\sigma/\Delta_\sigma$. The finite-index subgroup of $H_\sigma$ fixing components and half-edges has the cutting extension by $\Delta_\sigma$ with quotient $\prod_vP_{g_v,n_v}$. Full and procongruence completions agree in genus one and its genus-zero components. Boggi's stabilizer theorem and F.6's short-five-lemma argument identify this subgroup's completion with its closure. Passing through its finite-index extension does the same for $H_\sigma$ and gives (F.13). Multitwist independence and conjugacy separability give closedness of $H_\sigma,\Delta_\sigma$.

Here a transporter may be a finite union of cosets $H_{\sigma_t}r$, since different submulticurves of $\sigma_s$ may have the type of $\sigma_t$. Closedness of $H_{\sigma_t}$ supplies a finite quotient separating these finitely many cosets. Refine to it before the kernel argument in F.7. Equality modulo $V\Delta_{\sigma_s}$ then forces the same transporter component, since $\Delta_{\sigma_s}$ fixes all submulticurves. Hence $g'd^{-1}g^{-1}\in H_{\sigma_t}$ and the cofinality proof applies. This proves the transporter description and faithfulness, retaining graph symmetries.

Take an \etale atlas and its simplicial scheme of overlaps. Finite constructible sheaves satisfy descent. The analytic/\etale comparison on each overlap commutes with all face and degeneracy maps; its descent limit gives the stack comparison. This agrees with \citet[Proposition 13.9.17 and Corollary 13.9.19]{AF_BGH2020}. Galois pullback on the descent data supplies the coherent action and retains finite stabilizers.
\end{proof}

For $(1,1)$ the open group is $\widehat{\mathrm{SL}_2(\Z)}$, with central involution. Its boundary graph has one rational vertex, a labelled leg and a loop: exchanging the half-edges gives $A_\sigma=C_2$, while $P_{0,3}=1$. The boundary object's group is therefore $C_2$. For $(1,2)$ a rational tail carrying both markings leaves an elliptic component with one node, hence a factor $\widehat P_{1,1}$; a nonseparating node instead has a four-point rational normalization and an exchange of node branches. Formula (F.13) retains these stabilizers, including their $2$- and $3$-primary information.

\statement{F.16}{Open}{The precise remaining positive-genus automorphism comparison}
We have not established that the automorphism group of the two elliptic boundary categories, forgetting, and all boundary clutchings is $\GT$. The precise comparison asks whether restriction from the decorated procongruence outer tower E.13 to this diagram is an isomorphism, with identical oriented inertia and relabelling data. Injectivity requires detection at these levels; surjectivity requires extension of every such pair to the full decorated tower. We do not claim that existence of $\GT$ actions in positive genus remains open: E.12--E.13 give current results and their exact morphism collections.
\begin{proof}[Where our argument stops]
Formal completion and descent do not compute this automorphism group. Since $\widehat P_{1,1}$ has centre $C_2$, F.13's uniqueness-of-coherence argument needs additional work. Elliptic presentations also include nonseparating inertia and handle relations, whereas the genus-zero proof uses labelled splits and separating twists. E.10--E.11 concern specified open groups, and do not alone settle coherence and extension across all boundary bimodules. One must compare the precise decorated towers, retaining the central involution. The smallest diagram not settled here is exactly $(1,1)$--$(1,2)$. The label \emph{Open} records an unresolved comparison in this audit, not an assertion that no theorem exists elsewhere.
\end{proof}

\statement{F.17}{Known}{The rational operadic analogue and its exact limitation}
Over a characteristic-zero field $k$, define the rational analogue to be the parenthesized braid groupoids completed relative to their finite permutation quotients, with pure braid groups replaced by their Malcev prounipotent completions. Retain operadic composition, strand deletion and doubling, the braid symmetry, and the distinguished oriented inertia. Their structure-preserving automorphism group is $\mathrm{GT}(k)$. In the parenthesized chord-diagram realization it is $\mathrm{GRT}(k)$; an associator compares the realizations. The four- and five-point constraints are the symmetry equations and pentagon, and the identity-multiplier Lie algebra is the weight completion of $W=\mathfrak{grt}_1$. The exact sources are \citet[\S\S5--6, Propositions 5.1 and 5.5]{Drinfeld1990} and \citet[Proposition 3.4, Theorems 1--3, Corollary 4.1]{BarNatan1998}.
\begin{proof}[The corollary supplied by the cited formulation]
The algebraic completion is universal for homomorphisms from the pure groups to unipotent algebraic groups, while the discrete objects and finite permutation data are retained. The cited theorem computes automorphisms preserving the indicated operations, and identifies their equations on the two lowest nontrivial levels. It is the operadic, algebraically enriched analogue of the decorated tower automorphism calculation F.13. The change from braids to infinitesimal braids depends on an associator and is not canonical. Applying E.19--E.20 gives its motivic corollary: the mixed Tate group injects, and its equality with the full rational automorphism group follows conditionally from the all-weight upper bound in family 008. That last assertion remains explicitly \emph{claimed (OpenAI, 2026, unrefereed)} with the audit in E.26--E.28.

This does not prove that termwise Malcev completion of every Hom-set in the boundary category produces those parenthesized groupoids. A boundary quotient $G/\Delta_\sigma$ is generally a homogeneous space, not a torsor for $G$, and the exit category has noninvertible smoothing arrows. The parenthesized braid categories are groupoids. For instance the arrow from a boundary object to the open object of $C_4$ has no reverse arrow; completing its groups does not create one. Thus those bare categories cannot be identified. A further comparison of their \emph{structure-preserving automorphism groups} would require an algebraically enriched boundary construction and a reconstruction proof analogous to F.12; it is not imported silently from the groupoid theorem. This is the exact scope of the rational corollary supplied here. Locally constant coefficients in $D(\Q)$ do not by themselves implement Malcev completion, and neither rational construction yields $G_{\Q}=\GT$.
\end{proof}

\statement{F.18}{Open}{A stronger higher-completion comparison is not assumed}
One can also ask whether the nerve of the strict finite-category pro-object F.7, as a pro-object of stratified spaces, equals the full profinite stratified $\infty$-category $\operatorname{Gal}(X_{n,\C}/P_n)$. This requires comparison for all finite homotopy types and their constructible gluing, beyond finite sets. We do not establish it here. The first nontrivial test is $n=4$, whose boundary gluing retains a two-dimensional unstratified homotopy class even though its stratum groups are free.
\begin{proof}[Exact limitation]
BGH's Corollary 12.6.7 identifies the right-hand object with profinite stratified completion of the \emph{nerve} of $C_n$. F.2 instead takes strict quotient categories before nerves. Interchanging them requires cofinality or goodness for representations in all $\pi$-finite spaces. F.9's finite representables prove only the $0$-truncated coefficient comparison. All automorphism statements F.12--F.14 concern the strict completion defined here, whose finite-set comparison is proved. We leave the stronger statement unresolved without asserting a counterexample. No unchecked contractibility assertion from STRATA F.10 is used.
\end{proof}

